\documentclass[12pt]{article}
\usepackage[utf8]{inputenc}
\usepackage[T1]{fontenc}
\usepackage{hyperref}
\title{Teressa Language Resources for K-12 Curriculum Development}
\author{Compiled by Claude Sonnet 5 (claude-sonnet-5) for Andr\'as Kornai}
\date{}
\begin{document}
\maketitle

\section{Introduction}
Teressa (ISO 639-3 tef; Glottocode tere1272) is an endangered Nicobarese (Austroasiatic) language of the Nicobar Islands, India. Its native name is Luro. Searches found no standalone Teressa dictionary, grammar, primer or textbook. The material below is sparse.

\section{Primary Educational Resources}
No primers, textbooks or school programmes for Teressa turned up in the searches.

\section{Linguistic and Reference Materials}
\begin{itemize}
\item Teressa language, Wikipedia overview: \url{https://en.wikipedia.org/wiki/Teressa_language}
\item Nicobarese languages, Wikipedia overview: \url{https://en.wikipedia.org/wiki/Nicobarese_languages}
\item Numeral data for Teressa, collected by Vysakh R. and Anvita Abbi (Max Planck): \url{https://lingweb.eva.mpg.de/channumerals/Teressa.htm}
\item ASJP wordlist for Teressa: \url{https://asjp.clld.org/languages/TERESSA}
\item de Roepstorff (1875), a historical vocabulary that includes Teressa data (Glottolog record): \url{https://glottolog.org/resource/reference/id/111601}
\item Man's \emph{Dictionary and Grammar of the Nicobarese Language}. It covers Central Nicobarese, not Teressa specifically (bookseller listing): \url{https://vedamsbooks.com/books/dictionary_and_grammar_of_the_nicobarese_language}
\end{itemize}

\section{Religious and Community Resources}
None were found.

\section{Digital and Audio Resources}
\begin{itemize}
\item About Nicobarese, a site on the Nicobarese languages. The search summary said it includes fieldwork-based material on Luro (Teressa) and Sanenyo (Chowra) by Rahul Pachori and Vysakh R. I did not open the page to confirm this: \url{https://andamanese.org/nicobarese}
\end{itemize}

\section{Curriculum Development Context}
The language is not known to be taught in schools. Curriculum work would need new documentation, starting with the fieldwork materials above and the speakers themselves.

\end{document}